\chapter{Conclusion}
\label{ch:conclusion}

This project studied whether ML-based malicious-bitstream inspection can move into the FPGA deployment path for serverless FPGA systems. The target setting is a Coyote-style environment where user partial bitstreams are checked before partial reconfiguration. The goal is to add a hardware-backed security gate that can reject suspicious bitstreams before they enter the reconfigurable region.

We built a partial-bitstream inspection pipeline around three components. First, we generated a dataset with benign Coyote application bitstreams and RO-based malicious bitstreams. Second, we transformed each bitstream into a fixed-size byte-image tensor and trained CNN classifiers for binary detection. Third, we synthesized selected CNNs with hls4ml and evaluated their detection quality, latency, and FPGA resource overhead.

The results show that FPGA-resident ML inspection is feasible as a defense-in-depth mechanism. The 512$\times$7 W8A8 P50 manual-v1 model gives stronger accuracy-oriented metrics, including higher F1, lower false positive rate, higher ROC AUC, and perfect detection on the extended RO-footprint dataset. The 256$\times$7 W8A8 P50 manual-v1 model is the preferred deployment candidate because it gives a better hardware tradeoff. It achieves useful detection quality with 3.35 ms added latency and much lower BRAM overhead than the 512$\times$7 alternative.

The main takeaway is that deployment-oriented bitstream inspection must be evaluated across both ML and systems metrics. Software accuracy alone is not enough. A practical checker must also fit beside the FPGA shell and user logic, add limited latency to the PR path, and integrate cleanly into the deployment flow. This prototype shows a practical path toward such a checker, while leaving broader attack coverage, harder datasets, cross-platform validation, and full Coyote integration as future work.
